\section{Notes on the formalization}
\label{sec:formalization-notes}

This chapter records the wrong turns the formalization of
Katoh--Tanigawa's Theorem~5.5~\cite{katohTanigawa2011} took on the way
to the argument the earlier chapters present: statements whose formal
strength diverged from the source's, decompositions that did not
survive contact with the assembly they were meant to feed, and
obligations left with no artifact tracking them. Each section below is
one failure-mode class; within a class, an episode narrates what a
check did and did not read at the time a divergence was introduced, and
the check that eventually caught it. Because the subject here is the
formalization itself, Lean declarations, their types, and short code
excerpts appear directly in the prose below, rather than only as
parenthetical citations, as they do in the rest of this document.

Three patterns recur, and each episode below is one of them. A
\emph{genuine mis-factoring} is an internal decomposition --- an
abstraction layer, a splice of two arguments, a case split --- that did
not match the structure the final assembly needed, so the decomposition
was replaced without changing the underlying mathematics. A
\emph{source-faithfulness correction} is a Lean statement whose formal
strength diverged from Katoh--Tanigawa's --- weaker, stronger, or
differently quantified --- found and corrected without invalidating any
proof already built on it. A \emph{process failure} is a decision that
was sound when made but left with no artifact tracking it, or an error
that no check running at the time could have read.

Commits are cited by an eight-character abbreviated SHA linking to the
corresponding commit on GitHub; every cited commit postdates the
project's 2026-05-13 move to a standalone repository.

\subsection{Statement-faithfulness episodes}
\label{sec:formalization-notes-statement-faithfulness}

The episodes below are source-faithfulness corrections in the sense of
the introduction: a Lean statement's formal strength diverged from
Katoh--Tanigawa's, and the divergence was found and repaired without
invalidating proofs already built against it.

\subsubsection{A realization predicate weaker than its prose}

KT Theorem~5.5 asserts that a minimal $k$-dof graph admits a
panel-hinge realization whose rigidity matrix attains the maximal rank.
In KT's definition (KT~\S5.1) each hinge is a line: a nonzero
$2$-extensor contained in both endpoint panels. The formalization
introduced the realization predicate in June 2026
(commit~\href{https://github.com/bryangingechen/CombinatorialRigidity/commit/1cd26ccee6d8412a55740f2d9661934b5456701a}{\texttt{1cd26cce}})
as
\begin{alltt}
def HasFullRankRealization (k : \(\N\)) (G : Graph \(\alpha\) \(\beta\)) : Prop :=
  \(\exists\) Q : PanelHingeFramework k \(\alpha\) \(\beta\), Q.graph = G \(\wedge\)
    Q.toBodyHinge.RankHypothesis 0
\end{alltt}
with no condition on the hinge extensors. The omission makes the
predicate satisfiable for every connected graph. The framework
assigning every hinge the degenerate endpoint pair $(a_0, a_0)$ has all
support extensors zero; the constraint at an edge $uv$, that the screw
difference $S_u - S_v$ lie in the span of the edge's extensor, then
reads $S_u = S_v$, so on a connected graph every infinitesimal motion is
constant and the rank is the full $D(|V|-1)$. The realization conjunct
of the formalized Theorem~5.5 was satisfied by this witness
independently of the theorem's content.

The predicate had been introduced as rank bookkeeping for the induction
of KT~\S5 and was not compared against KT's definition of realization,
either at its introduction or at the one later redesign that reworked
it (the two-motive split of
\href{https://github.com/bryangingechen/CombinatorialRigidity/commit/e4693d6170908c480e4f06924f2ea0a95ce9ba3f}{\texttt{e4693d61}}'s
design pass re-read KT on the simple-versus-non-simple axis and did not
pose the question). The blueprint did not expose the divergence: the
definition's dependency-graph node was marked formalized beside prose
at the paper's strength, and none of the per-commit checks compared a
definition's Lean body with its prose --- the name-resolution check
verifies that pinned declarations exist, and the hypothesis audit reads
a theorem's hypotheses, of which a definition has none.

The divergence was found by a pre-build audit of the theorem's
hypothesis feeds
(commit~\href{https://github.com/bryangingechen/CombinatorialRigidity/commit/e4693d6170908c480e4f06924f2ea0a95ce9ba3f}{\texttt{e4693d61}},
2026-06-11), before any construction was built against the degenerate
witness. It was a weakness of statement selection rather than of proof:
the general-position motive introduced at the two-motive split is at
KT's strength, and the rank arguments beneath the bare conjunct
required no rework. The repair was scheduled at the motive redesign
that the same audit had already forced, so the definition surface was
reshaped once rather than twice, and landed as
(\href{https://github.com/bryangingechen/CombinatorialRigidity/commit/652ea99ffe115bd932f7b9da3324473cd8001d0f}{\texttt{652ea99f}})
\begin{alltt}
def HasPanelRealization (k n : \(\N\)) (G : Graph \(\alpha\) \(\beta\)) : Prop :=
  \(\exists\) (F : BodyHingeFramework k \(\alpha\) \(\beta\)) (normal : \(\alpha\) \(\to\) Fin (k + 2) \(\to\) \(\R\)),
    F.graph = G \(\wedge\)
    (\(\forall\) v \(\in\) V(G), normal v \(\neq\) 0) \(\wedge\)
    (\(\forall\) e u v, G.IsLink e u v \(\to\) F.supportExtensor e \(\neq\) 0 \(\wedge\)
      ExtensorInPanel (F.supportExtensor e) (normal u) \(\wedge\)
      ExtensorInPanel (F.supportExtensor e) (normal v)) \(\wedge\) \dots
\end{alltt}
whose nonzero-normal and nonzero-extensor conjuncts are the
nondegeneracy conditions of KT's definition. The check this produced is
the cheapest-witness audit (\texttt{blueprint/CLAUDE.md}, the
definition-faithfulness gate): when a definition modeling a paper
notion is marked formalized, its blueprint prose is read against the
cheapest satisfying witness of the Lean body, not against the paper.
